% LaTeX companion of 01-tensor-products-and-matrix-multiplication.typ. Both formats are editable.
\chapter{tensor product and matrix multiplication}

这里我们放弃陈述两个 over 同一 field 的 vector spaces 的 tensor product 的 algebraic definition，直接看应用的。(完整的 definition 是两个 over 同一个 ring 的 modules \(A,B\)，取它们 free abelian group generated by \(A \times B\)，再 quotient 掉一个用来形成 bilinearity 的 subgroup，就是它们的 tensor product。当这两个东西是 vector spaces 时，它们并且是 isomorphic to 其对应的 bilinear functional vector space的。)

\begin{definition}{(not rigorous) \kn{Tensor product}}

令 \(V,W\) be vector spaces over \(\mathbb{F}\)，我们定义一个 vector 与 vector 之间的 tensor product operation：

\(\otimes :\left( {v,w} \right)\rightarrow v \otimes w\)

s.t. 对于 \(V,W\) 的 basis \(e_{1},\ldots,e_{n}\) 和 \(\varepsilon_{1},\ldots,\varepsilon_{m}\)，我们给每个 \(\left( {e_{i},\varepsilon_{j}} \right)\) 都赋予一个不同的 image \(e_{i} \otimes \varepsilon_{j}\)，其 over \(\mathbb{F}\) 具有 bilinear 性。

对于

\(V \otimes W := \text{span}\left( \left\{ {e_{i} \otimes e_{j}:1 \leq i \leq n,1 \leq j \leq m} \right\} \right)\)

\(V \otimes W\) 也是一个 over \(\mathbb{F}\) 的 vector space，且 \(\text{dim}\left( {V \otimes W} \right) = mn\)。

\end{definition}

\begin{proposition}{representing tensor product as matrix}

如果我们有两个矩阵：

\(A = \begin{pmatrix}
a_{11} & a_{12} \\
a_{21} & a_{22}
\end{pmatrix},\quad B = \begin{pmatrix}
b_{11} & b_{12} \\
b_{21} & b_{22}
\end{pmatrix}\)

可以把它们的 tensor product \(A \otimes B\) 表示为：

\(A \otimes B = \begin{pmatrix}
{a_{11}B} & {a_{12}B} \\
{a_{21}B} & {a_{22}B}
\end{pmatrix} = \begin{pmatrix}
{a_{11}b_{11}} & {a_{11}b_{12}} & {a_{12}b_{11}} & {a_{12}b_{12}} \\
{a_{11}b_{21}} & {a_{11}b_{22}} & {a_{12}b_{21}} & {a_{12}b_{22}} \\
{a_{21}b_{11}} & {a_{21}b_{12}} & {a_{22}b_{21}} & {a_{22}b_{22}} \\
{a_{21}b_{21}} & {a_{21}b_{22}} & {a_{22}b_{21}} & {a_{22}b_{22}}
\end{pmatrix}\)

can verify：这个表示是符合 tensor product 的 bilinearity 的。

(我们可以把 \(A\)，\(B\) 分别看作 \(m \times n\)，\(p \times q\) dim 的向量，这个 \(A \otimes B\) 是 \(m \times n \times p \times q\) dim 的向量。)

\end{proposition}

\begin{definition}{outer product of two vectors}

对于 \(v \in \mathbb{F}^{n}\)，\(w \in \mathbb{F}^{m}\)，我们定义它们的 outer product \(v \otimes w\)为：

\(vw := v \otimes w^{T} = \begin{pmatrix}
v_{1} \\
\ldots \\
v_{n}
\end{pmatrix} \otimes \begin{pmatrix}
w_{1} & \ldots & w_{m}
\end{pmatrix} = \begin{pmatrix}
{v_{1}w^{T}} \\
\ldots \\
{v_{n}w^{T}}
\end{pmatrix} = \begin{pmatrix}
{w_{1}v} & \ldots & w_{m}
\end{pmatrix}\)

\end{definition}

\section{matrix product through outer product}

\begin{theorem}{representing matrix multiplication by outer products}

对于 \(m \times n\) 的矩阵 \(A\) 和 \(n \times k\) 的矩阵 \(B\)，we have：

\(AB = \sum_{i = 1}^{n}A_{\ast i} \otimes B_{i \ast}\)

\end{theorem}

\begin{proof}

In md.

\end{proof}

\begin{qlremark}{Remark}

两个 vectors 的 outer product，一\textbf{定是一个 rank 为 1 的 matrix}。因为它每一 column 都是 \(v\) 的一个倍数。

但是我们可以累加 \(n\) 个这样的 outer products，从而得到一个可以是 full rank \(n\) 的矩阵。(我们知道当然，矩阵相加不保留 rank。)

\end{qlremark}
